\section{Introduction}
\label{sec:introduction}

\PARstart{L}{arge} language model (LLM) agents increasingly answer analytics questions by writing and running SQL on
a user's behalf. Such an agent usually connects with the user's own credential, so it can read everything the user
can read. Consider an analyst at a school district who asks an agent: \emph{``How many charter schools in Alameda
County offer a free-meal program?''} The answer needs three columns of two tables and the rows of one county. The
analyst's entitlement, however, covers every school, every test score and every contact field in the database. If the
agent processes untrusted content, for example a note in a retrieved document or a text value stored in the
database, an injected instruction can make it read and repeat data the question never needed~\cite{greshake2023not,debenedetti2024agentdojo}.
The agent acts as a confused deputy~\cite{hardy1988confused}: it holds the user's authority for one narrow purpose but
can be steered to use all of it.

Least privilege says a subject should run with only the privileges its job needs~\cite{saltzer1975protection,dennis1966programming,nist80053r5}.
For agents, recent systems apply this to \emph{tools}: they decide per task which tools, and sometimes which tool
arguments, are allowed~\cite{zhu2025miniscope,shi2025progent,cai2025agentsentry,zhu2026igac,zheng2026intentcap,ruan2024toolemu}.
For a data agent, however, the relevant tool is typically a single \texttt{run\_query} interface. Once it is allowed,
the data the agent can reach is unchanged: still the full entitlement. Databases have long supported finer control
through views, column grants and row-level security~\cite{stonebraker1975query,rizvi2004extending,pgrls,oraclevpd},
and access-controlled text-to-SQL systems enforce a \emph{user's} policy~\cite{fei2026rbac,klisura2026role,miyamoto2026pccsql,kim2026securemcp}.
Practitioners write static scoped views per role. None of these narrows the data an agent can reach to what the
\emph{current task} needs.

We propose \emph{task-scoped entitlements} (TSE). Before an agent starts, a trusted broker infers a relational scope
from the natural-language task: which tables, columns and rows it needs. The broker compiles that scope into an
ephemeral PostgreSQL role and \texttt{security\_barrier} views that are owned by the user's role, hands the agent a
login for that role as its only credential, and drops everything when the task ends. Because the views are owned by
the user's role, the database bounds them by the user's entitlement even when the inference is wrong. When the
inferred scope is insufficient, the agent can file a structured widening request, which is bounded by the
entitlement, decided by a policy, and logged.

\subsection{Contributions}
\begin{itemize}
\item[\textbf{C1}] A formal model of task scopes: entitlement, task scope, needed data, sufficiency, minimality, a
column/row/cell over-exposure ratio (OER), and a widening bound, with a design guarantee stated under explicit trust
assumptions (Section~\ref{sec:model}).
\item[\textbf{C2}] A scope-inference procedure that combines an LLM proposal with an automatically built semantic
layer, deterministic schema linking, grounding of row filters in the task text, foreign-key closure, and a
conservativeness parameter $k$ (Section~\ref{sec:inference}).
\item[\textbf{C3}] Database enforcement independent of the model: compilation to ephemeral roles and
\texttt{security\_barrier} views owned by the user role, with idempotent teardown (Section~\ref{sec:enforcement}).
\item[\textbf{C4}] A controlled widening protocol that turns inference errors into bounded, auditable requests
(Section~\ref{sec:widening}).
\item[\textbf{C5}] A preregistered evaluation on \nTasks{} public text-to-SQL tasks (BIRD and Spider) with frontier
and open-weight agents in two agent frameworks, four baselines, indirect prompt injection, widening abuse, and
\bypassAttempts{} enforcement-bypass attempts (Sections~\ref{sec:setup}--\ref{sec:results}).
\end{itemize}

\subsection{Summary of findings}
Under synthetic-policy entitlements, the default TSE scope reduces the columns an agent can reach by a factor of
\oerColReduction{} (geometric mean) and the reachable cells by a factor of \oerCellReduction{} relative to the full
entitlement. The scope is sufficient for the gold query on \suffTse\% of tasks; this inference quality is the
load-bearing assumption of the approach, and we quantify it directly. With widening, execution accuracy changed by
\dExHaikuTse{} percentage points for Claude Haiku~5.5 (95\% CI [\dExLoHaikuTse, \dExHiHaikuTse]) and by
\dExSonnetTse{} points for Claude Sonnet~5.5 (95\% CI [\dExLoSonnetTse, \dExHiSonnetTse]) relative to full access, a smaller loss than static role views (\dExHaikuBone{} points) or prompt-level schema
pruning with full credentials (\dExHaikuBthree{} points) for Haiku~5.5. A more conservative scope ($k{=}2$) brought
accuracy to \ablExKtwo\% against \ablExBzero\% for B0 on an exploratory subset.
Under indirect prompt injection that targets out-of-task data, exfiltration fell from \exfSonnetCtxBzero\% (B0) to
\exfSonnetCtxTse\% (TSE) for Sonnet~5.5 and from \exfHaikuCtxBzero\% to \exfHaikuCtxTse\% for Haiku~5.5; nearly all remaining leaks used
approved widening, and with widening denied exfiltration was \exfSonnetCtxDeny\% and \exfHaikuCtxDeny\%. None of
our \bypassAttempts{} scripted bypass attempts reached out-of-scope data. Provisioning takes \medCompile{}\,ms
(median) and teardown \medTeardown{}\,ms.

\subsection{What we do not claim}
TSE reduces the data reachable for a task; it does not prevent all leakage. It does not protect against inference or
aggregation over data inside the scope, covert or timing channels, exfiltration of in-scope data, a malicious
authorized user, or a compromised broker, and table and column names remain visible in the catalog. Its benefit
depends on two assumptions that can fail: the quality of scope inference (too narrow causes failures or widening;
too broad reduces the benefit) and the correctness of database enforcement. Views, row-level security and query
modification are established techniques that we use unchanged; the new parts are listed in
Section~\ref{sec:whatsnew}. We make no claim about production readiness: the schemas are public benchmarks and the
entitlements are synthetic.

The rest of the paper is organized as follows. Section~\ref{sec:threat} states the threat model,
Section~\ref{sec:related} positions the work, Sections~\ref{sec:model}--\ref{sec:enforcement} present the model and
system, Sections~\ref{sec:setup}--\ref{sec:results} the evaluation, and Sections~\ref{sec:discussion}--\ref{sec:conclusion}
discuss threats, limitations and future work.
